\chapter{Conclusion}

\section{Kết luận tổng thể HW02}

Bốn phạm vi A-FR03, B-FR09, C-FR14 và D-M09 đã có artifact theo workflow
Domain Testing: đặc tả/oracle, domain model, EP/BVA, test case, automation hoặc
execution evidence, defect/gap analysis, traceability và feature closure.
Tổng cộng 177 test case được thiết kế; 97 case có verdict gồm 55 PASS và 42
FAIL; 80 case BLOCKED; 0 NOT RUN.

Kết quả không chứng minh hệ thống hoạt động đúng. Pool A còn 29 BLOCKED và tám
FAIL; Pool B còn lỗi Critical về percent formula, trusted total và identity;
Pool C còn authorization/data-integrity defect; Pool D tái hiện năm lỗi Backend
dùng chung, thêm hai lỗi Mobile và còn 20 BLOCKED.

\section{Kết luận riêng cho Mobile}

D-M09 hoàn tất workflow kỹ thuật: 50 designed; 30 executed; 19 PASS; 11 FAIL;
20 BLOCKED; 0 NOT RUN. Baseline được thu trên Expo Web/Chromium tại viewport
Mobile \(440 \times 956\). Chín screenshot bổ sung trên một thiết bị iOS thật
chứng minh kết nối LAN và một số trạng thái UI/keyboard, nhưng khác baseline và
không phải native automation. Vì vậy chưa có căn cứ kết luận Android, rotation,
screen reader, native semantic/focus hoặc hiệu năng hoạt động đúng.

Trạng thái chất lượng D-M09 là \texttt{NOT ACCEPTED}. Cần sửa các defect về
formula, trusted pricing, authentication/identity và Mobile UI; chốt oracle cho
20 case bị chặn; duy trì API endpoint theo environment; sau đó chạy regression
trên Expo Web và native iOS/Android với oracle, metadata và trace phù hợp.

\section{Giới hạn của kết luận}

“Workflow technically complete” chỉ mô tả tính đầy đủ và truy vết được của
artifact trong phạm vi hiện tại. Nó không đồng nghĩa mọi requirement đã được
thực thi, không đồng nghĩa coverage 100\%, và không phải quyết định chấp nhận
SUT. Human sign-off, xử lý gap và regression vẫn là điều kiện trước khi chốt
chất lượng cuối cùng.
